\chapter{Decision Trees}
\label{ch:trees}
Trees are the most interview-friendly model: you can compute a split by hand in two minutes, draw
the result, and explain it to a non-technical manager. They are also the building block of the
models that win most tabular competitions (\Cref{ch:ensembles}), so their quirks --- greediness,
instability, bias towards many-valued features --- matter far beyond this chapter.

\section{Impurity}
\prototype{A node with 8 candidates who accepted an offer and 6 who declined: how ``mixed'' is it?}

A classification tree splits so that child nodes are purer than the parent. With class proportions
$p_1,\dots,p_K$ in a node:
\begin{definition}
	\begin{gather*}
		\text{entropy } \Ent=-\sum_kp_k\log_2p_k,\qquad
		\text{Gini } G=1-\sum_kp_k^2=\sum_kp_k(1-p_k),\\
		\text{misclassification error } 1-\max_kp_k.
	\end{gather*}
\end{definition}
All three are $0$ for a pure node and maximal for a uniform one (binary: entropy $1$ bit, Gini
$0.5$). Gini is the probability that two labels drawn at random from the node disagree; entropy is
the expected number of bits to encode a label. They almost always choose the same split; Gini avoids
a logarithm and is scikit-learn's default.

\begin{definition}
	The \vocab{information gain} of a split of node $S$ into children $S_v$ is
	\[ \IG=\Ent(S)-\sum_v\frac{\abs{S_v}}{\abs S}\Ent(S_v), \]
	the drop from parent impurity to the size-weighted child impurity. The same construction with
	Gini gives the \emph{Gini decrease}.
\end{definition}

\begin{remark}
	\trap Misclassification error is a poor splitting criterion: it is not strictly concave, so many
	splits that clearly improve purity leave it unchanged (e.g.\ $(40,40)\to(30,10)+(10,30)$ and
	$(40,40)\to(20,40)+(20,0)$ both have error $0.25$, but the second creates a pure node). Use it for
	pruning, not growing.
\end{remark}

\section{A split by hand}
\prototype{Which of CTC, city and role best predicts whether a candidate accepts an offer?}

\begin{center}
	\small
	\begin{tabular}{rllll}
		\toprule
		\# & CTC & City & Role & Accepted? \\ \midrule
		\gv{c09_trees/table}
		\bottomrule
	\end{tabular}
\end{center}

\begin{example}[Information gain of each feature]
	Parent: $\gv{c09_trees/npos}$ yes, $\gv{c09_trees/nneg}$ no, so
	$\Ent=-\frac8{14}\log_2\frac8{14}-\frac6{14}\log_2\frac6{14}=\gv{c09_trees/Hp}$ and Gini
	$=1-(8/14)^2-(6/14)^2=\gv{c09_trees/Gp}$.
	\begin{itemize}
		\ii \textbf{CTC} ($\gv{c09_trees/CTC_parts}$): child entropies $\gv{c09_trees/CTC_High_H}$ and
		$\gv{c09_trees/CTC_Low_H}$, weighted $\frac7{14}\cdot0.5917+\frac7{14}\cdot0.8631=
			\gv{c09_trees/CTC_condH}$, so $\IG=\gv{c09_trees/CTC_IG}$.
		\ii \textbf{City} ($\gv{c09_trees/City_parts}$): weighted child entropy $\gv{c09_trees/City_condH}$,
		$\IG=\gv{c09_trees/City_IG}$.
		\ii \textbf{Role} ($\gv{c09_trees/Role_parts}$): weighted child entropy $\gv{c09_trees/Role_condH}$,
		$\IG=\gv{c09_trees/Role_IG}$.
	\end{itemize}
	The root splits on \textbf{\gv{c09_trees/best}}. Gini decreases rank the same way
	($\gv{c09_trees/CTC_Gdec}$, $\gv{c09_trees/City_Gdec}$, $\gv{c09_trees/Role_Gdec}$). The parent
	Gini and the CTC gain match scikit-learn's impurity values for a depth-1 tree.
\end{example}

\begin{center}
	\begin{tikzpicture}[
			node/.style={draw, rounded corners, align=center, font=\small, fill=blue!5, minimum width=2.6cm},
			leaf/.style={draw, rounded corners, align=center, font=\small, fill=green!8, minimum width=2.6cm},
			level distance=1.6cm, sibling distance=5cm, edge from parent/.style={draw, -Stealth}]
		\node[node] {all 14\\ 8 yes / 6 no\\ $\Ent=\gv{c09_trees/Hp}$}
		child {node[leaf] {CTC = High\\ 6 yes / 1 no\\ $\Ent=\gv{c09_trees/CTC_High_H}$} edge from parent node[left, font=\scriptsize] {High~}}
		child {node[leaf] {CTC = Low\\ 2 yes / 5 no\\ $\Ent=\gv{c09_trees/CTC_Low_H}$} edge from parent node[right, font=\scriptsize] {~Low}};
	\end{tikzpicture}
\end{center}

\subsection{The bias towards many-valued features: gain ratio}
City and Role have \emph{identical} information gain, but Role splits three ways. Taken to the
extreme, a candidate-ID feature with 14 distinct values gives 14 pure leaves and the maximal gain
$\gv{c09_trees/id_IG}$ while being useless for prediction. C4.5 corrects this with the
\vocab{gain ratio} $\IG/\mathrm{SplitInfo}$, where
$\mathrm{SplitInfo}=-\sum_v\frac{\abs{S_v}}{\abs S}\log_2\frac{\abs{S_v}}{\abs S}$ is the entropy of the
split itself. SplitInfo is $\gv{c09_trees/City_SI}$ for City, $\gv{c09_trees/Role_SI}$ for Role and
$\log_214=\gv{c09_trees/id_SI}$ for the ID, giving gain ratios $\gv{c09_trees/City_GR}$,
$\gv{c09_trees/Role_GR}$ and $\gv{c09_trees/id_GR}$. (Gain ratio over-corrects towards
near-trivial splits with tiny SplitInfo, so C4.5 only considers splits with above-average gain.)

\begin{remark}[ID3, C4.5, CART]
	\vocab{ID3}: categorical features, multiway splits, information gain, no pruning.
	\vocab{C4.5}: adds continuous thresholds, gain ratio, missing-value handling and error-based
	pruning. \vocab{CART} (what scikit-learn implements): strictly \emph{binary} splits, Gini for
	classification and squared error for regression, cost-complexity pruning. A binary tree can
	still express a multiway categorical split as a sequence of binary ones.
\end{remark}

\section{Growing a tree}
\prototype{Eight values of a continuous feature: only seven thresholds are worth trying.}

\subsection{Continuous features}
Sort the values; only midpoints between consecutive \emph{distinct} values can change the
partition, so a feature with $n$ distinct values has at most $n-1$ candidate thresholds (and only
those between differently-labelled neighbours can be optimal). With pre-sorting, scanning all
thresholds for one feature costs $O(n)$ by updating class counts incrementally.
\begin{example}[Threshold scan]
	$x=(2,3,4.5,5,6.5,7,8,9.5)$ with labels $(0,0,0,1,0,1,1,1)$; parent Gini $\gv{c09_trees/thr_parent}$.
	\begin{center}
		\small
		\begin{tabular}{cccc}
			\toprule
			threshold & left (positives/size) & right & weighted Gini \\ \midrule
			\gv{c09_trees/thr_rows}
			\bottomrule
		\end{tabular}
	\end{center}
	Two thresholds tie at $\gv{c09_trees/thr_best_g}$; scikit-learn's stump picks
	$x\le\gv{c09_trees/thr_best}$, the first.
\end{example}

\subsection{Categorical features and missing values}
A categorical feature with $q$ levels admits $2^{q-1}-1$ binary partitions. For binary
classification (and regression) there is a shortcut: order the levels by their positive rate (mean
target) and only consider the $q-1$ splits in that order --- optimal for Gini and squared error.
That ordering is a form of target encoding, with the same leakage risk if done on the full data
(\Cref{ch:dataprep}). Libraries that take only numbers need one-hot or ordinal encoding first;
LightGBM and CatBoost handle categories natively.

Missing values: C4.5 sends a fractional example down every branch; CART uses \vocab{surrogate
	splits} (another feature mimicking the primary split); XGBoost and LightGBM learn a
\emph{default direction} per split; recent scikit-learn trees accept \texttt{NaN} and learn which
side missing values go to. With 10\% of entries set to \texttt{NaN} in both train and test, a depth-5
tree on the breast-cancer data still reaches test accuracy $\gv{c09_trees/nan_acc}$.

\subsection{Stopping}
Recursion stops when a node is pure, or hits \code{max_depth}, \code{min_samples_split},
\code{min_samples_leaf}, \code{max_leaf_nodes}, or a minimum impurity decrease. Greedy
growth is not optimal: finding the smallest consistent tree is NP-hard, and a split with zero gain
(XOR!) may be needed to enable great splits below it.

\section{Regression trees}
\prototype{Six points that jump from about 5 to about 21: split where the jump is.}
Each leaf predicts the mean of its targets; the split criterion is the reduction in the sum of
squared errors (equivalently, variance reduction).
\begin{example}[One regression split]
	$x=1,\dots,6$, $y=(5,6,5,20,22,21)$; total SSE $\gv{c09_trees/r_sst}$. The best threshold is
	$x\le\gv{c09_trees/r_t}$ with leaf means $\gv{c09_trees/r_l}$ and $\gv{c09_trees/r_r}$ and remaining
	SSE $\gv{c09_trees/r_sse}$ --- identical to \texttt{DecisionTreeRegressor(max\_depth=1)}.
\end{example}
A regression tree is a piecewise-constant function: it cannot extrapolate beyond the range of the
training targets and approximates a straight line by a staircase. That is why trees forecast
trending time series badly unless you model differences, and why ``trees for regression'' answers
should mention it.

\section{Pruning}
\prototype{An unpruned tree reaches 100\% training accuracy; its test accuracy peaked several levels
	earlier.}

Grown to purity, a tree memorises noise. \vocab{Pre-pruning} stops early (depth, leaf size);
\vocab{post-pruning} grows fully and cuts back. CART's \vocab{cost-complexity pruning} minimises
\[ R_\alpha(T)=R(T)+\alpha\abs T \]
(training error plus $\alpha$ per leaf). As $\alpha$ grows, the subtree whose removal increases error
least per leaf removed (the ``weakest link'') is collapsed first, giving a nested sequence of
subtrees; pick $\alpha$ by cross-validation. Note the parallel with $E+\lambda\Omega$ in
\Cref{ch:biasvar}.

\begin{example}[Breast-cancer data]
	The unpruned tree has depth $\gv{c09_trees/full_depth}$, $\gv{c09_trees/full_leaves}$ leaves and
	test accuracy $\gv{c09_trees/full_test}$. Along the pruning path, the best test accuracy
	$\gv{c09_trees/ccp_best_test}$ is reached at $\alpha=\gv{c09_trees/ccp_best_alpha}$ with
	$\gv{c09_trees/ccp_best_leaves}$ leaves.
\end{example}

\begin{figure}[ht]
	\centering
	\begin{tikzpicture}
		\begin{groupplot}[group style={group size=2 by 1, horizontal sep=1.4cm},
				width=0.48\linewidth, height=0.38\linewidth, grid=major, grid style={gray!20},
				axis lines=left, ticklabel style={font=\scriptsize},
				legend style={draw=none, fill=none, font=\scriptsize}]
			\nextgroupplot[xlabel={max depth}, ylabel={accuracy}, legend pos=south east, ymin=0.88, ymax=1.01]
			\addplot[blue!80!black, thick, mark=*, mark size=1pt] table[x=d,y=train] {gen/ml/data/c09_trees-depth.dat};
			\addplot[red!80!black, thick, mark=square*, mark size=1pt] table[x=d,y=test] {gen/ml/data/c09_trees-depth.dat};
			\legend{train, test}
			\nextgroupplot[xlabel={$\alpha$ (log)}, xmode=log, ylabel={leaves}, axis y line*=left]
			\addplot[ForestGreen, thick, mark=*, mark size=1pt, const plot] table[x=alpha,y=leaves] {gen/ml/data/c09_trees-ccp.dat};
		\end{groupplot}
	\end{tikzpicture}
	\caption{Left: training accuracy rises to 1 with depth while test accuracy plateaus. Right:
		cost-complexity pruning removes leaves as $\alpha$ grows.}
	\label{fig:tree-depth}
\end{figure}

\section{What trees are good for}
\prototype{OA: ``Which of the following can decision trees be used for?''}

\textbf{Strengths.} Classification and regression (and multi-output); mixed numeric and categorical
features; no scaling needed --- predictions are invariant to any monotone transform of a feature
(log-transforming every feature leaves the breast-cancer tree's test predictions identical, checked
in the script); automatic interactions (a path is a conjunction of conditions); missing-value
strategies; interpretable rules; built-in feature-importance scores; fast prediction ($O(\text{depth})$);
feature selection and discretisation as by-products; base learners for ensembles.

\textbf{Weaknesses.} High variance: resampling the training set changes the root split --- over 200
bootstrap samples of the breast-cancer data the depth-1 tree chose $\gv{c09_trees/root_distinct}$
different root features, the commonest only $\gv{c09_trees/root_top_share}$ of the time. Axis-aligned
boundaries approximate a diagonal by a staircase (\Cref{fig:tree-regions}). No extrapolation. Greedy
and myopic. Biased towards features with many split points. Probability estimates from small leaves
are coarse.

\begin{figure}[ht]
	\centering
	\begin{tikzpicture}
		\begin{groupplot}[group style={group size=2 by 1, horizontal sep=0.8cm},
				width=0.5\linewidth, height=0.4\linewidth, xmin=-1.5, xmax=2.5, ymin=-1.1, ymax=1.6,
				ticklabel style={font=\scriptsize}, title style={font=\small}, axis on top]
			\nextgroupplot[title={depth 2 (\gv{c09_trees/moon_leaves2} leaves)}]
			\addplot[scatter, only marks, mark=square*, mark size=1.9pt, point meta=\thisrow{c},
				colormap={rb}{color=(blue!12) color=(red!12)}, scatter/use mapped color={draw=mapped color, fill=mapped color}]
			table[x=x,y=y] {gen/ml/data/c09_trees-region2.dat};
			\addplot[only marks, mark=o, mark size=1pt, blue!70!black] table[x=x,y=y] {gen/ml/data/c09_trees-m0.dat};
			\addplot[only marks, mark=+, mark size=1.3pt, red!80!black] table[x=x,y=y] {gen/ml/data/c09_trees-m1.dat};
			\nextgroupplot[title={depth 12 (\gv{c09_trees/moon_leaves12} leaves)}, yticklabels={}]
			\addplot[scatter, only marks, mark=square*, mark size=1.9pt, point meta=\thisrow{c},
				colormap={rb}{color=(blue!12) color=(red!12)}, scatter/use mapped color={draw=mapped color, fill=mapped color}]
			table[x=x,y=y] {gen/ml/data/c09_trees-region12.dat};
			\addplot[only marks, mark=o, mark size=1pt, blue!70!black] table[x=x,y=y] {gen/ml/data/c09_trees-m0.dat};
			\addplot[only marks, mark=+, mark size=1.3pt, red!80!black] table[x=x,y=y] {gen/ml/data/c09_trees-m1.dat};
		\end{groupplot}
	\end{tikzpicture}
	\caption{Decision regions are unions of axis-aligned rectangles. A deep tree carves small boxes
		around individual noisy points.}
	\label{fig:tree-regions}
\end{figure}

\begin{remark}
	\trap ``Decision trees cannot handle categorical data'' (they can in principle; scikit-learn's
	implementation needs encoding). ``Decision trees need normalised features'' (they do not).
	``Decision trees are linear classifiers'' (no: piecewise-constant, axis-aligned). ``A deeper tree
	always generalises better'' (no: training accuracy rises, test accuracy plateaus or falls).
\end{remark}

\begin{moral}
	Split greedily on the largest impurity drop, stop or prune before the leaves memorise noise,
	and expect high variance --- which is exactly what ensembles fix.
\end{moral}

\section\problemhead

\begin{problem}[Entropy and Gini]
	\onechili
	A node holds 30 samples: 10 of class A, 10 of B, 10 of C. Compute its entropy (bits) and Gini.
	It is split into $(10,10,0)$ and $(0,0,10)$; compute the information gain.
	\begin{hint}
		Uniform over three classes: $\log_23$.
	\end{hint}
	\begin{sol}
		Entropy $\log_23=\gv{c09_problems/p9a_H}$, Gini $1-3\cdot\frac19=\gv{c09_problems/p9a_G}$.
		Children: $(10,10,0)$ has entropy $1$, $(0,0,10)$ has $0$; weighted
		$\frac{20}{30}\cdot1=0.667$; $\IG=\gv{c09_problems/p9a_IG}$.
	\end{sol}
\end{problem}

\begin{problem}[Best split]
	\twochili
	12 emails, 6 spam. Feature A splits them into $(5\text{ spam},1\text{ ham})$ and $(1,5)$. Feature
	B splits them into $(3,0)$ and $(3,6)$. Which split does information gain prefer? Which does Gini
	prefer?
	\begin{hint}
		Compute both weighted child impurities.
	\end{hint}
	\begin{sol}
		A: each child has entropy $\gv{c09_problems/p9b_hA}$, IG $=\gv{c09_problems/p9b_igA}$; Gini
		decrease $\gv{c09_problems/p9b_gA}$. B: children entropies $0$ and $\gv{c09_problems/p9b_hB}$,
		weighted $\frac9{12}\cdot0.918$, IG $=\gv{c09_problems/p9b_igB}$; Gini decrease
		$\gv{c09_problems/p9b_gB}$. Both prefer A, although B creates a pure node.
	\end{sol}
\end{problem}

\begin{problem}[Regression split]
	\onechili
	$x=(1,2,3,4)$, $y=(2,4,10,12)$. Find the best single split by squared error and the two leaf
	predictions.
	\begin{hint}
		Try $x\le1.5$, $x\le2.5$, $x\le3.5$.
	\end{hint}
	\begin{sol}
		SSE for the three thresholds: $\gv{c09_problems/p9c_s1}$, $\gv{c09_problems/p9c_s2}$,
		$\gv{c09_problems/p9c_s3}$. Best is $x\le2.5$ with leaves $3$ and $11$.
	\end{sol}
\end{problem}

\begin{problem}[True or false]
	\onechili
	(a) A fully grown tree on data without duplicate inputs has zero training error. (b) Standardising
	features changes a decision tree's predictions. (c) Trees can capture feature interactions without
	explicit product features. (d) A regression tree can predict a value larger than every training
	target. (e) Gini impurity of a binary node is at most 0.5.
	\begin{hint}
		Think about leaves and monotone transforms.
	\end{hint}
	\begin{sol}
		(a) true. (b) false. (c) true. (d) false: leaves average training targets. (e) true:
		$2p(1-p)\le\frac12$.
	\end{sol}
\end{problem}
